% concatenated from Weak_action_representability_of_2-nilpotent_groups.tex
\documentclass[reqno,hidelinks,a4paper]{amsart}
\usepackage[T1]{fontenc}
\usepackage{latexsym, amssymb, amsmath}
\usepackage{enumitem}
\usepackage{amsthm}
\usepackage{tikz-cd}
\usepackage{orcidlink}
\usepackage{doi}
\renewcommand{\doi}[1]{DOI:~\href{https://doi.org/#1}{#1}}
\usepackage{url}
\usepackage[utf8]{inputenc}
\usepackage[british]{babel}
\usepackage[all]{xy}
\usepackage{hyperref}
\usepackage{calrsfs}
\usepackage{mathabx}
\usepackage{xcolor}
\usepackage{blkarray}
\usepackage{microtype}
\usepackage[textsize=scriptsize]{todonotes}
\usepackage{thmtools,mathtools}
\usepackage{nameref,hyperref,cleveref}

\theoremstyle{plain}
\newtheorem{theorem}[subsection]{Theorem}
\newtheorem{proposition}[subsection]{Proposition}
\newtheorem{corollary}[subsection]{Corollary}	
\newtheorem{lemma}[subsection]{Lemma}

\theoremstyle{definition}
\newtheorem{definition}[subsection]{Definition}

\theoremstyle{remark}
\newtheorem{remark}[subsection]{Remark}
\newtheorem{example}[subsection]{Example}
\newtheorem{examples}[subsection]{Examples}
\newtheorem{problem}[subsection]{Open problem}

\renewcommand{\labelitemi}{\(\cdot\)}
\setcounter{tocdepth}{1}
\numberwithin{equation}{section}

\newenvironment{tfae}
{\renewcommand{\labelenumi}{(\roman{enumi})}
\begin{enumerate}}
{\end{enumerate}}

\newcommand{\noproof}{\hfil\qed}

\newcommand*{\CondName}[1]{{\rm (#1)}}
\newcommand*{\LACC}{\CondName{LACC}}

\newcommand{\C}{\ensuremath{\mathcal{C}}}
\newcommand{\F}{\mathbb{F}}         
\newcommand{\E}{\mathcal{E}}
\newcommand{\Z}{\mathbb{Z}}

\DeclareMathOperator{\bchar}{char}
\DeclareMathOperator{\bZ}{Z}
\DeclareMathOperator{\Ker}{Ker}
\DeclareMathOperator{\Aut}{Aut}
\DeclareMathOperator{\Act}{Act}
\DeclareMathOperator{\DAct}{DAct}
\DeclareMathOperator{\Hom}{Hom}
\DeclareMathOperator{\SplExt}{SplExt}
\DeclareMathOperator{\Imm}{Im}
\DeclareMathOperator{\Der}{Der}
\DeclareMathOperator{\Inn}{Inn}
\DeclareMathOperator{\M}{M}
\DeclareMathOperator{\USGA}{USGA}
\DeclareMathOperator{\id}{id}

\newcommand{\Grp}{\ensuremath{\mathsf{Grp}}}
\newcommand{\Ab}{\ensuremath{\mathsf{Ab}}}
\newcommand{\Set}{\ensuremath{\mathsf{Set}}}
\newcommand{\Assoc}{\ensuremath{\mathsf{Assoc}}}
\newcommand{\Lie}{\ensuremath{\mathsf{Lie}}}
\newcommand{\Leib}{\ensuremath{\mathsf{Leib}}}
\newcommand{\CAssoc}{\ensuremath{\mathsf{CAssoc}}}
\newcommand{\AbAlg}{\ensuremath{\mathsf{AbAlg}}}
\newcommand{\Nil}{\ensuremath{\mathsf{Nil}}}
\newcommand{\SplExtX}{\ensuremath{\mathsf{SplExt}(X)}}

\newdir{>}{{}*:(1,-.2)@^{>}*:(1,+.2)@_{>}}
\newdir{<}{{}*:(1,+.2)@^{<}*:(1,-.2)@_{<}}
\newdir{>>}{{}*!/3.5pt/:(1,-.2)@^{>}*!/3.5pt/:(1,+.2)@_{>}*!/7pt/:(1,-.2)@^{>}*!/7pt/:(1,+.2)@_{>}}
\newdir{ >}{{}*!/-8pt/@{>}}

\date{}

\begin{document}

\title[Weak action representability of \texorpdfstring{\(2\)}{2}-nilpotent groups]{Weak action representability\\of \texorpdfstring{\(2\)}{2}-nilpotent groups}

\author[A.~Dioguardi Burgio]{Alessandro Dioguardi Burgio~\orcidlink{0000-0003-0856-0144}}
%\author[M.~Mancini]{Manuel Mancini\(^{*}\)~\orcidlink{0000-0003-2142-6193}}
\author[M.~Mancini]{Manuel Mancini~\orcidlink{0000-0003-2142-6193}}
\author[T.~Van~der Linden]{Tim Van~der Linden~\orcidlink{0000-0001-5474-5356}}

\email{alessandro.dioguardiburgio@unipa.it}
\email{manuel.mancini@unipa.it; manuel.mancini@uclouvain.be}
\email{tim.vanderlinden@uclouvain.be}

\address[Alessandro Dioguardi Burgio]{Dipartimento di Ingegneria, Università degli Studi di Palermo, Viale delle Scienze, 90128 Palermo, Italy.}

\address[Manuel Mancini]{Dipartimento di Matematica e Informatica, Università degli Studi di Palermo, Via Archirafi 34, 90123 Palermo, Italy.}

\address[Manuel Mancini, Tim Van der Linden]{Institut de Recherche en Mathématique et Physique, Université catholique de Louvain, chemin du cyclotron 2 bte L7.01.02, B--1348 Louvain-la-Neuve, Belgium.}

\address[Tim Van der Linden]{Mathematics \& Data Science, Vrije Universiteit Brussel, Pleinlaan 2, B--1050 Brussel, Belgium}

%\thanks{\(^*\)Corresponding Author.}

\begin{abstract}
    In this article, we investigate the representability of actions of the category \(\Nil_2(\Grp)\) of \(2\)-nilpotent groups. We first provide an algebraic characterisation of derived actions in \(\Nil_2(\Grp)\) by determining a universal strict general actor of an object \(X\), which turns out to be the group \(\Aut_c(X)\) of central automorphisms of \(X\). We also characterise the morphisms \(B \to \Aut_c(X)\) that define an action of \(B\) on \(X\) in \(\Nil_2(\Grp)\).

    We then show that \(\Nil_2(\Grp)\) is not action representable, and that the existence of a weak representation is related to the amalgamation property. Using the construction of an amalgam of a suitable family of abelian subgroups of \(\Aut_c(X)\), we prove that the category \(\Nil_2(\Grp)\) is weakly action representable, and that a weak representing object can be chosen to be an abelian group.

    Finally, we show that \(\Nil_2(\Grp)\) is not locally algebraically cartesian closed.
\end{abstract}

\subjclass[2020]{08A35; 08C05; 16W22; 17A36; 18E13; 20F18; 20F28}
\keywords{Action accessible category; action representable  category; amalgamation property; central automorphism; local algebraic cartesian closedness; nilpotent group; split extension}

\maketitle

\section*{Introduction}\label{Introduction}
\emph{Internal actions} were defined in~\cite{BJK} by F.~Borceux, G.~Janelidze and G.~M.~Kelly in order to extend the correspondence between actions and split extensions from the category of groups to arbitrary semi-abelian categories~\cite{Semi-Ab}. In some of those categories, internal actions are exceptionally well behaved, in the sense that the actions on each object~\(X\) are \emph{representable}. More explicitly, this means that there exists an object~\([X]\) such that the functor
\[
    \Act(-,X)\cong \SplExt(-,X),
\]
sending an object~\(B\) to the set of actions of~\(B\) on~\(X\) (equivalently, to the set of isomorphism classes of split extensions of~\(B\) by~\(X\)), is naturally isomorphic to the functor \(\Hom(-,[X])\).

The notion of action representability in a semi-abelian category was further investigated in~\cite{BJK2}, where it is shown there that the category of commutative associative algebras over a field is not action representable. Subsequently, it was proved in~\cite{Tim} that, over an infinite field of characteristic different from~\(2\), the only action representable varieties of non-associative algebras are the category \(\Lie\) of Lie algebras and the category \(\AbAlg\) of abelian algebras. The strength of this notion naturally motivated the introduction of weaker related concepts.

In~\cite{act_accessible}, D.~Bourn and G.~Janelidze introduced the notion of an \emph{action accessible} category so as to encompass relevant examples that do not fit into the frame of action representable categories, such as rings, associative algebras, and Leibniz algebras~\cite{loday}. It was later proved by A.~Montoli in~\cite{Montoli} that every \emph{Orzech category of interest}~\cite{Orzech} is action accessible. Moreover, in~\cite{Casas}, the authors showed that every Orzech category of interest admits a weaker form of representing object, called the \emph{universal strict general actor} (USGA).

More recently, G.~Janelidze introduced in~\cite{WRA} the notion of a \emph{weakly action representable} category. Rather than requiring that, for every object \(X\) of a semi-abelian category~\(\C\), there exist an object \([X]\) and a natural isomorphism
\[
    \Act(-,X)\cong\Hom_{\C}(-,[X]),
\]
one only asks for the existence of an object \(T=T(X)\) of \(\C\) together with a natural monomorphism of functors
\[
    \Act(-,X)\rightarrowtail \Hom_{\C}(-,T).
\]
Such an object \(T\) is called a \emph{weak representing object} of \(X\). Examples of weakly action representable categories include the category of associative algebras~\cite{WRA} and the category of Leibniz algebras~\cite{CigoliManciniMetere}.

The notion of weak action representability was further investigated in~\cite{WRAAlg} in the setting of \emph{operadic} varieties of non-associative algebras over a field. In that work, the authors exhibited new examples of weakly action representable categories, including the variety of commutative associative algebras and the varieties of \(2\)-nilpotent (commutative and anti-commutative) algebras.

Another important observation made by  G.~Janelidze in~\cite{WRA} is that every weakly action representable category is action accessible. The converse, however, does not hold. Indeed, J.~R.~A.~Gray observed in~\cite[Proposition~2.2]{Gray} that weak action representability is closely related to the existence of a so-called \emph{amalgam}~\cite{kiss}, a notion that had already appeared in~\cite{BJK2} in connection with action representability, and he proved that the varieties of \(n\)-solvable groups are not weakly action representable, for any \(n \geq 3\). This result was later extended in~\cite{XabiMancini} to the varieties of \(2\)-solvable groups and of \(k\)-nilpotent groups, for any \(k \geq 3\). There, the authors also proved that the varieties of \(n\)-solvable Lie algebras (\(n \geq 2\)) and \(k\)-nilpotent Lie algebras (\(k \geq 3\)) fail to be weakly action representable. These are the first examples of action accessible varieties of non-associative algebras which are not weakly action representable.

The main goal of this article is to investigate the representability of actions in the category \(\Nil_2(\Grp)\) of \(2\)-nilpotent groups. After recalling the necessary background in \Cref{SecPrel}, we give an algebraic characterisation of (external) actions in \(\Nil_2(\Grp)\). More precisely, we characterise the actions of a \(2\)-nilpotent group \(B\) on another \(2\)-nilpotent group~\(X\) such that the corresponding semidirect product \(B \ltimes X\) is still \(2\)-nilpotent.

As a consequence, we obtain an explicit description of the universal strict general actor of~\(X\), which turns out to be the group \(\Aut_c(X)\) of \emph{central automorphisms} of~\(X\). We also provide a complete characterisation of the \emph{acting morphism}, that is, of those group homomorphisms \(B \to \Aut_c(X)\) that define an action of \(B\) on \(X\) in the category \(\Nil_2(\Grp)\).

In \Cref{sec_rep}, we investigate weak action representability in the category \(\Nil_2(\Grp)\). We first prove in \Cref{thm_ar} that \(\Nil_2(\Grp)\) is not action representable. Then, similarly to what was done in~\cite{WRAAlg} for the variety of commutative associative algebras, we show that, for a \(2\)-nilpotent group \(X\), the existence of a weak representation for the actions on \(X\) is closely related to the so-called \emph{amalgamation property} (AP)~\cite{kiss}. Finally, by constructing the amalgam of a suitable family of abelian subgroups of \(\Aut_c(X)\), we prove that \(\Nil_2(\Grp)\) is weakly action representable, and that as a weak representing object of \(X\) we may choose an abelian group. This result is analogous to the case of \(2\)-nilpotent Lie algebras. Indeed, it was shown in~\cite[Theorem 2.21]{WRAAlg} that the variety of \(2\)-nilpotent Lie algebras over a field~\(\F\) with \(\bchar(\F)\neq 2\) is weakly action representable, with a weak representing object of an object \(X\) being an abelian algebra.

Finally, in \Cref{sec_LACC} we study \emph{local algebraic cartesian closedness} \LACC{} for the category \(\Nil_2(\Grp)\), a notion introduced by J.~R.~A.~Gray in his Ph.D.\ thesis~\cite{GrayThesis} in connection with the concept of \emph{algebraic exponentiation} (see~\cite{BournGray,GrayACS}). In the semi-abelian setting, this property provides a natural categorical framework for the study of commutators and internal actions. In the context of varieties of non-associative algebras over a field, it has been shown that the condition \LACC{} essentially characterises the variety of Lie algebras (see~\cite{GM-VdL2,GM-VdL3,GrayLie}), so that, in this context, it is equivalent to action representability. At present, it remains an open problem whether local algebraic cartesian closedness implies (weak) action representability, or whether a counterexample exists. Here, using the characterisation of the condition \LACC{} asking that the comparison map
\[
    (B \flat X) + (B \flat Y) \to B \flat (X+Y),
\]
where \(B \flat X\) denotes the kernel of the morphism \([\id_B,0]\colon {B+X\to B}\), is an isomorphism, we prove that the category \(\Nil_2(\Grp)\) is not \LACC{}. This suggests a further line of investigation into the \LACC{} property for other subvarieties of \(\Grp\).

\section{Preliminaries}\label{SecPrel}

The notion of \emph{semi-abelian} category was introduced in~\cite{Semi-Ab} by G.~Janelidze, L.~Márki and W.~Tholen in order to provide a categorical setting which would capture categorical-algebraic properties of non-abelian algebraic structures, such as groups, rings and Lie algebras. We recall that a category~\(\C\) is
semi-abelian when it is pointed, Barr-exact~\cite{Barr}, Bourn-protomodular~\cite{Bourn} and has finite coproducts.

A notion that can be explored in the context of semi-abelian categories is that of split extension.

\begin{definition}
    Let \(\C\) be a semi-abelian category and let \(B\), \(X\) be objects of \(\C\). A \emph{split extension} of \(B\) by \(X\) is a diagram
    \begin{equation*}\label{eq:split_ext}
        \begin{tikzcd}
            X\arrow [r, "k"]
            &A \arrow[r, shift left, "p"] &
            B\ar[l, shift left, "s"]
        \end{tikzcd}
    \end{equation*}
    in \(\C\) such that \(p \circ s = \id_B\) and \((X,k)\) is a kernel of \(p\).
\end{definition}

For any object \(X\) of \(\C\), one can define the functor
\[
    \SplExt(-,X) \colon \C^{op} \to \Set
\]
which maps any object \(B\) of \(\C\) to the set \(\SplExt(B,X)\) of isomorphism classes of split extensions of \(B\) on \(X\), and any arrow \(f\colon B'\to B\) to the \emph{change of base} morphism \(f^*\colon \SplExt(B,X) \to \SplExt(B',X)\) given by pulling back along \(f\).

A feature of semi-abelian categories is that one can define a notion of \emph{internal action}. Internal actions correspond to split extensions via a semidirect product construction; it turns out that for our purposes we need no explicit description of what is an internal action, since split extensions are easier to work with, especially in categories of groups with operations. We refer the interested reader to~\cite{BJK2}, where the equivalence between the two concepts is described in detail.

For us here, it is sufficient to observe that if we fix an
object \(X\) of \(\C\), internal actions on \(X\) give rise to a functor
\[
    \Act(-,X)\colon \C^{op} \to \Set
\]
and a natural isomorphism \(\Act(-,X)\cong \SplExt(-,X)\) (see~\cite{BJK2}). This justifies the terminology in the definition that follows.

\begin{definition}\cite{BJK2}
    A semi-abelian category \(\C\) is said to be \emph{action representable} if for any object \(X\) in \(\C\), the functor \(\SplExt(-,X)\) is representable. This means that there exists an object \([X]\) of \(\C\) and a natural isomorphism of functors
    \[
        \SplExt(-,X) \cong \Hom_{\C}(-,[X]).
    \]
\end{definition}

The prototype examples of action representable categories are the category \(\Grp\) of groups and the category \(\Lie\) of Lie algebras over a commutative unital ring. In the first case, every action of \(B\) on \(X\) is represented by a group homomorphism \(B \to \Aut(X)\), where \(\Aut(X)\) is the group of automorphisms of \(X\). In the second case, any action of \(B\) on \(X\) is represented by a Lie algebra homomorphism \(B \to \Der(X)\), where \(\Der(X)\) is the Lie algebra of derivations of \(X\).

However, the notion of an action representable category has proved to be quite restrictive: for instance, the authors of~\cite{Tim} showed that, over an infinite field \(\F\) with \(\bchar(\F)\neq 2\), the only action representable varieties of non-associative algebras~\cite{VdL-NAA} are the variety of abelian algebras and the variety of Lie algebras.

\begin{remark}\label{BJ}
    We recall that a semi-abelian category \(\C\) is action representable if and only if, for every object \(X\) of \(\C\), the category \(\SplExtX\) of split extensions in~\(\C\) with kernel \(X\) admits a terminal object of the form
    \[
        \begin{tikzcd}
            X \arrow[r]
            & {[X] \ltimes X} \arrow[r, shift left] &
            {[X]\text{.}} \arrow[l, shift left]
        \end{tikzcd}
    \]
    This condition can be relaxed by requiring only that every object of \(\SplExtX\) is \emph{accessible}, i.e., it admits a morphism into a \emph{subterminal} (or \emph{faithful}) object. Here, we recall that an object is faithful if there exists at most one morphism into it. This gives rise to the following notion, introduced by D.~Bourn and G.~Janelidze in~\cite{act_accessible} in order to compute centralisers of normal subobjects and equivalence relations.
\end{remark}

\begin{definition}\cite{act_accessible}
    A semi-abelian category \(\C\) is \emph{action accessible} if, for any object \(X\) of \(\C\), every object in \(\SplExtX\) is accessible.
\end{definition}

The concept of action accessibility makes it possible to include a broader class of categories that fail to be action representable, such as the category of rings and the variety of associative algebras.

We observe that, since the existence of a terminal object in \(\SplExtX\) is stronger than every object being accessible, it follows from \Cref{BJ} that
\[
    \emph{action representability}\Rightarrow\emph{action accessibility.}
\]

A different way to weaken action representability is to relax the requirements on the representing object~\([X]\), rather than the existence of a terminal object in \(\SplExtX\). This is the approach adopted in~\cite{Casas}, where it was shown that every \emph{Orzech category of interest}~\cite{Orzech} admits a so-called \emph{universal strict general actor} (USGA). It was then proved by A.~Montoli in~\cite{Montoli} that every Orzech category of interest is action accessible. Moreover, it was shown in~\cite{CigoliManciniMetere} that, for every object~\(X\) of an Orzech category of interest \(\C\), there exists a natural monomorphism of functors
\[
    \tau \colon \SplExt(-,X) \rightarrowtail \Hom_{\C'}(U(-),\USGA(X)),
\]
where \(\C'\) is a category containing \(\C\) as a full subcategory, and \(U \colon \C \to \C'\) is the forgetful functor.

More recently, G.~Janelidze introduced in~\cite{WRA} the notion of weak action representability.

\begin{definition}\cite{WRA}
    A semi-abelian category \(\C\) is said to be \emph{weakly action representable} if for every object \(X\) in \(\C\), the functor \(\SplExt(-,X)\) admits a weak representation. This means that there exists an object \(T=T(X)\) of \(\C\) and a natural monomorphism of functors
    \[
        \tau\colon \SplExt(-,X) \rightarrowtail \Hom_{\C}(-,T).
    \]
    A morphism \((\varphi\colon B \to T) \in \Imm(\tau_B)\) is called \emph{acting morphism}.
\end{definition}

Notice that every action representable category is weakly action representable. In this case, \(T=[X]\) is the actor of \(X\), \(\tau\) is a natural isomorphism and every arrow \(\varphi\colon B \to [X]\) is an acting morphism. Other examples of weakly action representable categories are the variety \(\Assoc\) of associative algebras~\cite{WRA}, the variety \(\Leib\) of Leibniz algebras~\cite{CigoliManciniMetere}, and the varieties of \(2\)-nilpotent (commutative and anti-commutative) algebras (see~\cite{WRAAlg}). More generally, if in an Orzech category of interest \(\C\) each \(\USGA(X)\) is an object of the category, then \(\C\) is weakly action
representable (see~\cite[Corollary 4.2]{CigoliManciniMetere}).

\begin{remark}
    We observe that in the above mentioned examples, the weakly representing object \(T=T(X)\) is quite easy to describe. However, this is not always the case. For instance, the variety \(\CAssoc\) of commutative associative algebras is also weakly action representable, but the algebra that could be expected to be a weakly representing object is not commutative. Nevertheless, in~\cite{WRAAlg} a large colimit, along with the amalgamation property is used to naturally construct a weakly representing object.
\end{remark}

Another important observation made by G.~Janelidze in~\cite{WRA} is that every weakly action representable category is action accessible. We thus have that
\[
    \textit{action representability}\Rightarrow\textit{weak action representability}\Rightarrow\textit{action accessibility.}
\]
The author ended with an open question: whether reasonably mild conditions may be found on a semi-abelian category under which the second implication may be reversed.

This question was addressed by J.~R.~A.\ Gray, who showed in~\cite{Gray} that the notion of weak action representability is related to the so-called \emph{amalgamation property} (AP)~\cite{kiss}. Using (AP), he established sufficient conditions under which a Birkhoff subcategory of an action representable category fails to be weakly action representable. In particular, he proved that the varieties of \(n\)-solvable groups are not weakly action representable for any \(n \geq 3\). This result was later adapted in~\cite{XabiMancini} to prove that the varieties of \(k\)-nilpotent groups (\(k \geq 3\)), as well as the variety of \(2\)-solvable groups, are not weakly action representable. In the same paper, it was also shown that the varieties of \(k\)-nilpotent Lie algebras (\(k \geq 3\)) and the varieties of \(n\)-solvable Lie algebras (\(n \geq 2\)) fail to be weakly action representable. Note, however, that this is no longer true for \(k=2\): indeed, it was proved in~\cite[Theorem~2.21]{WRAAlg} that the variety of \(2\)-nilpotent anti-commutative algebras over a field \(\F\) with \(\bchar(\F)\neq 2\) is weakly action representable, with a weak representing object on an object \(X\) being an abelian algebra.

In this article, we show that an analogue of Theorem~2.21 of~\cite{WRAAlg} holds for the variety \(\Nil_2(\Grp)\) of \(2\)-nilpotent groups. More precisely, in the next section we provide an algebraic characterisation of actions in the variety of \(2\)-nilpotent groups, and we give an explicit description of the universal strict general actor of a \(2\)-nilpotent group \(X\). We then use this description in \Cref{sec_rep} to prove that \(\Nil_2(\Grp)\) is weakly action representable, and that a weak representing object of \(X\) is an abelian group.

\section{Actions and central automorphisms of \texorpdfstring{\(2\)}{2}-nilpotent groups}\label{sec_act}

In this section we aim to describe actions in the variety of \(2\)-nilpotent groups. We start by recalling the following definition.

\begin{definition}
    Let \(X\) be a group. The \emph{lower central series} of \(X\) is defined as
    \[
        X^{0}=X, \quad X^{k}=[X^{k-1},X],
    \]
    for any \(k \in \mathbb{N}\).

    The group \(X\) is said to be \emph{nilpotent of class \(k\)} if \(X^{k-1}\neq \{1\}\) and \(X^{k}=\{1\}\).
\end{definition}

We denote by \(\Nil_k(\Grp)\) the subvariety of \(\Grp\) consisting of all groups of nilpotency class at most~\(k\).

\begin{remark}
    Note that the class of all nilpotent groups (i.e., for any nilpotency class) does not form a variety; however, it does when the nilpotency class is bounded. For this reason, in the rest of the paper, when we say that a group \(X\) is \(2\)-nilpotent, we mean that \(X\) is an object of \(\Nil_2(\Grp)\), i.e., \(X\) is nilpotent of class at most \(2\). More precisely, the category \(\Nil_2(\Grp)\) consists of all groups \(X\) such that the \emph{commutator subgroup}
    \[
        [X,X]=\langle [x,y] \mid x,y \in X \rangle
    \]
    lies in the center \(\bZ(X)\) of \(X\). Here, as usual,
    \[
        [x,y]=xyx^{-1}y^{-1}
    \]
    denotes the commutator of the elements \(x\) and \(y\).
\end{remark}

Since the variety \(\Nil_2(\Grp)\) is an Orzech category of interest~\cite{Orzech}, it is convenient to describe internal actions in terms of the so-called \emph{derived actions}.

\begin{definition}
    Let
    \begin{equation}\label{diag:Spl}
        \begin{tikzcd}
            X\arrow [r, "k"]
            &A \arrow[r, shift left, "p"] &
            B\ar[l, shift left, "s"]
        \end{tikzcd}
    \end{equation}
    be a split extension of \(2\)-nilpotent groups. The map
    \[
        \xi \colon B \times X \to X \colon (b,x) \mapsto b \ast x
    \]
    defined by
    \[
        b \ast x = s(b)k(x)s(b)^{-1},
    \]
    is called the \emph{derived action} of \(B\) on \(X\) associated with \eqref{diag:Spl}.
\end{definition}

In fact, for every Orzech category of interest there is a natural isomorphism of functors
\[
    \Act(-,X) \cong \DAct(-,X),
\]
where \(\DAct(-,X)\) denotes the functor of derived actions on a fixed object \(X\). We refer the reader to~\cite[Remark 3.16 and Proposition 3.17]{Mancini_Thesis} for a detailed proof.

Given a map
\[
    \xi \colon B \times X \to X \colon (b,x) \mapsto b \ast x,
\]
one can always define a binary operation on the cartesian product \(B \times X\) by
\[
    (b,x) \cdot (b',x')=(bb',x (b \ast x')).
\]
Via~\cite[Theorem 2.4]{Orzech}, this defines a \(2\)-nilpotent group structure on \(B \times X\) if and only if \(\xi\) is a derived action of \(B\) on \(X\) in \(\Nil_2(\Grp)\). This, in turn, is equivalent to a set of conditions on the map \(\xi\), as explained in the following proposition, which is a special case of in~\cite[Proposition 1.1]{Datuashvili}.

\begin{proposition}\label{PropAct}
    \((B \times X,\cdot)\) is a \(2\)-nilpotent group if and only if the following conditions hold:
    \begin{enumerate}
        \item \(\xi\) is a derived action of \(B\) on \(X\) in the category \(\Grp\) of groups, that is:
              \begin{itemize}
                  \item[(i)] \(1 \ast x = x\);
                  \item[(ii)] \((bb') \ast x = b \ast (b' \ast x)\);
                  \item[(iii)] \(b \ast (xx') = (b \ast x)(b \ast x')\);
              \end{itemize}

        \item \((b \ast x)x^{-1} \in \bZ(X)\cap X^B\),
    \end{enumerate}
    for any \(b\), \(b' \in B\) and \(x\), \(x' \in X\), where
    \[
        X^B=\{x \in X \mid b\ast x=x, \; \forall b\in B\}
    \]
    is the set of fixed points of the action\footnote{If \(\xi\) is a derived action of \(B\) on \(X\) in \(\Grp\), then one may check that \(X^B\) is a subgroup of \(X\).}. The resulting \(2\)-nilpotent group is the \emph{semidirect product} of \(B\) and \(X\), denoted by \(B \ltimes X\).
\end{proposition}

\begin{proof}
    Suppose that \((B \times X,\cdot)\) is a \(2\)-nilpotent group. Then, \(\xi\) is a derived action in the category \(\Grp\). Moreover, for any \((b,x)\), \((b',x')\), \((c,y)\in B\times X\), the equality
    \[
        [(b,x),(b',x')] \cdot (c,y)=(c,y) \cdot [(b,x),(b',x')]
    \]
    gives
    \begin{equation}\label{eq_2nil}
        \begin{split}
            x(b \ast x')([b,b']b' \ast x^{-1}) & ([b,b'] \ast (x'^{-1}y))                                                                            \\
            =                                  & y\, c \ast \bigl(x(b \ast x')\bigl([b,b']b' \ast x^{-1}\bigr)\bigl([b,b'] \ast x'^{-1}\bigr)\bigr).
        \end{split}
    \end{equation}
    Now, for \(b'=c=1\) and \(x=1\), using the fact that \(b \ast 1=1\), we obtain
    \[
        (b \ast x')x'^{-1}y=y(b \ast x')x'^{-1},
    \]
    and therefore
    \[
        (b \ast x')x'^{-1}\in \bZ(X).
    \]
    Similarly, if we take \(b=1\) and \(x=y=1\), we get
    \[
        c \ast \bigl((b \ast x')x'^{-1}\bigr)=(b \ast x')x'^{-1},
    \]
    so that
    \[
        (b \ast x')x'^{-1}\in X^B.
    \]
    Conversely, suppose that conditions~(1) and~(2) hold. Since \(\xi\) is a derived action of~\(B\) on \(X\) in \(\Grp\), it follows that \((B \times X,\cdot)\) is a group. It remains to prove that it is \(2\)-nilpotent, that is,
    \[
        [(b,x),(b',x')] \cdot (c,y)=(c,y)\cdot [(b,x),(b',x')]
    \]
    for every \((b,x)\), \((b',x')\), \((c,y)\in B\times X\). Since both \(B\) and \(X\) are \(2\)-nilpotent, it is enough to verify that \eqref{eq_2nil} holds.

    Now, the condition \((b \ast x)x^{-1}\in X^B\) implies that
    \[
        (bb')\ast x=(b'b)\ast x
    \]
    for any \(b\), \(b'\in B\) and \(x\in X\). Indeed, one has
    \begin{align*}
        ((bb') \ast x)((b'b) \ast x)^{-1} = & ((bb') \ast x)((b'b) \ast x^{-1})                                                      \\
        =                                   & (b \ast (b' \ast x)) (b' \ast (b \ast x^{-1}))                                         \\
        =                                   & (b \ast ((b' \ast x) x^{-1})) (b \ast x) (b' \ast ((b \ast x^{-1})x))(b' \ast x^{-1} ) \\
        =                                   & (b' \ast x)x^{-1} (b \ast x) (b \ast x^{-1})x (b' \ast x^{-1} )                        \\
        =                                   & (b' \ast x)x^{-1} x (b' \ast x^{-1} )                                                  \\
        =                                   & (b' \ast x)(b' \ast x^{-1})=1.
    \end{align*}
    Therefore, \([b,b'] \ast x=x\) and \eqref{eq_2nil} reduces to
    \[
        gy = y(c\ast g),
    \]
    where
    \[
        g=x(b\ast x')(b\ast x^{-1})x'^{-1}.
    \]
    To conclude, observe that
    \[
        g=\bigl((b'\ast x)x^{-1}\bigr)^{-1}[\,b'\ast x,\; b\ast x'\,](b\ast x')x'^{-1}.
    \]
    Moreover, the commutator subgroup \([X,X]\) is fixed by the action of \(B\). Indeed, let \(b\in B\) and \(x\), \(x'\in X\). By condition~(2), there exist \(z_x\), \(z_{x'}\in \bZ(X)\) such that
    \[
        b\ast x = z_x x
        \quad\text{and}\quad
        b\ast x' = z_{x'}x'.
    \]
    It follows that
    \[
        b\ast [x,x']=[b\ast x,b\ast x']=[z_xx,z_{x'}x']=[x,x'],
    \]
    since \(z_x\) and \(z_{x'}\) are central.

    Now, since \(X^B\) is a subgroup of \(X\) and
    \[
        [X,X]\leq \bZ(X)\cap X^B,
    \]
    we deduce from the above expression for \(g\) that
    \[
        g\in \bZ(X)\cap X^B.
    \]
    Hence
    \[
        gy = yg = y(c\ast g),
    \]
    so that \eqref{eq_2nil} holds. Therefore, \((B \times X,\cdot)\) is a \(2\)-nilpotent group.
\end{proof}

\begin{remark}
    We observe that, for any split extension \eqref{diag:Spl} and the corresponding derived action
    \[
        \xi \colon B \times X \to X,
    \]
    there is an isomorphism of split extensions
    \begin{equation*}
        \begin{tikzcd}
            X \arrow[r, "\iota_2"] \arrow[d, equal]
            & B \ltimes X \arrow[r, shift left, "\pi_1"] \arrow[d, "\theta"]
            & B \arrow[l, shift left, "\iota_1"] \arrow[d, equal] \\
            X \arrow[r, "k"]
            & A \arrow[r, shift left, "p"]
            & B \arrow[l, shift left, "s"]
        \end{tikzcd}
    \end{equation*}
    where \(\iota_1\), \(\iota_2\), and \(\pi_1\) are the canonical injections and projection, and \(\theta\) is defined by
    \[
        \theta(b,x)=s(b)k(x),
    \]
    for every \((b,x)\in B\ltimes X\).
\end{remark}

\begin{remark}\label{rmk}
    Since \(\xi\) is an action of \(B\) on \(X\) in the category \(\Grp\), the assignment
    \[
        b \mapsto \varphi_b = b \ast -
    \]
    defines a group homomorphism \(\varphi \colon B \to \Aut(X)\). Moreover, the condition
    \[
        (b \ast x)x^{-1} \in \bZ(X)
    \]
    means that \(\varphi_b\) is a \emph{central automorphism} of \(X\):
    \[
        \varphi_b \in \Aut_c(X)=\{ f \in \Aut(X) \mid f(x)x^{-1} \in \bZ(X), \; \forall x \in X \}.
    \]
    It is straightforward to check that \(\Aut_c(X)\) is a normal subgroup of \(\Aut(X)\), but is not, in general, an object of \(\Nil_2(\Grp)\). For instance, if \(X=\Z_2\times \Z_2\), then
    \[
        \Aut_c(X)=\Aut(X)=\operatorname{GL}_2(\Z_2) \cong S_3
    \]
    which is not nilpotent.

    Conversely, let
    \[
        \varphi \colon B \to \Aut_c(X) \colon b \mapsto b \ast -
    \]
    be a group homomorphism such that
    \[
        b' \ast \bigl((b \ast x)x^{-1}\bigr)=(b \ast x)x^{-1}
    \]
    for any \(b\), \(b' \in B\) and \(x \in X\). Then, one can define a split extension
    \begin{equation*}
        \begin{tikzcd}
            X \arrow[r, "\iota_2"]
            & B \ltimes X \arrow[r, shift left, "p_1"]
            & B \arrow[l, shift left, "\iota_1"]
        \end{tikzcd}
    \end{equation*}
    in \(\Nil_2(\Grp)\), where the group structure on \(B \ltimes X\) is defined by
    \[
        (b,x)\cdot (b',x')=(bb',x(b\ast x')).
    \]
    However, a generic morphism \(B \to \Aut_c(X)\) need not give rise to a split extension in \(\Nil_2(\Grp)\), as the following example shows.
\end{remark}

\begin{example}
    Consider the \(2\)-nilpotent group
    \[
        B = \langle x, a, b \mid [x, a]=b^{-1}, [x, b] = x^5 = a^5 = b^5=1 \rangle.
    \]
    Let \(X=\Z_5 \times \Z_5 \times \Z_5\) and let \(\varphi \colon B \rightarrowtail \Aut(X)\) be the group monomorphism defined by
    \[
        \varphi(x) =
        \begin{pmatrix}
            1 & 1 & 0 \\
            0 & 1 & 0 \\
            0 & 0 & 1
        \end{pmatrix}, \quad
        \varphi(a) =
        \begin{pmatrix}
            1 & 0 & 0 \\
            0 & 1 & 1 \\
            0 & 0 & 1
        \end{pmatrix}, \quad
        \varphi(b)=
        \begin{pmatrix}
            1 & 0 & -1 \\
            0 & 1 & 0  \\
            0 & 0 & 1
        \end{pmatrix}.
    \]
    It is immediate to see that the relations are preserved.

    Let \(\xi\) be the action of \(B\) on \(X\) induced by \(\varphi\), that is,
    \[
        g \ast \underline{v}=\varphi(g) \underline{v}, \quad \forall g \in B, \; \forall \underline{v}=\begin{pmatrix}
            v_1 \\ v_2 \\ v_3
        \end{pmatrix} \in X.
    \]
    Since \(X\) is abelian, every automorphism of \(X\) is central. However, \(\xi\) is not an action in the variety \(\Nil_2(\Grp)\), since
    \[
        a \ast \begin{pmatrix}
            0 \\ 1 \\ 1
        \end{pmatrix} + \begin{pmatrix}
            -0 \\ -1 \\ -1
        \end{pmatrix} = \begin{pmatrix}
            1 & 0 & 0 \\
            0 & 1 & 1 \\
            0 & 0 & 1
        \end{pmatrix} \begin{pmatrix}
            0 \\ 1 \\ 1
        \end{pmatrix} + \begin{pmatrix}
            0 \\ -1 \\ -1
        \end{pmatrix} = \begin{pmatrix}
            0 \\ 1 \\ 0
        \end{pmatrix}
    \]
    and
    \[
        x* \Bigg( a \ast \begin{pmatrix}
                0 \\ 1 \\ 1
            \end{pmatrix} + \begin{pmatrix}
                0 \\ -1 \\ -1
            \end{pmatrix} \Bigg)=\begin{pmatrix}
            1 & 1 & 0 \\
            0 & 1 & 0 \\
            0 & 0 & 1
        \end{pmatrix}\begin{pmatrix}
            0 \\ 1 \\ 0
        \end{pmatrix}=\begin{pmatrix}
            1 \\ 1 \\ 0
        \end{pmatrix} \neq \begin{pmatrix}
            0 \\ 1 \\ 0
        \end{pmatrix}.
    \]
    In fact, it is shown in~\cite{XabiMancini} that the corresponding semidirect product \(B\ltimes X\) has nilpotency class \(3\).
\end{example}

We can now claim the following result.

\begin{theorem}\label{main_thm}
    Let \(X\) be a \(2\)-nilpotent group and let \(U \colon \Nil_2(\Grp) \to \Grp\) be the forgetful functor.
    \begin{itemize}
        \item[(i)] There exists a natural monomorphism of functors
              \[
                  \rho \colon \SplExt(-,X) \rightarrowtail \Hom_\Grp(U(-),\Aut_c(X)).
              \]

        \item[(ii)] A group homomorphism
              \[
                  \varphi \colon U(B) \to \Aut_c(X) \colon b \mapsto b \ast -
              \]
              belongs to \(\Imm(\rho_B)\) if and only if
              \begin{equation}\label{act_mor}
                  b' \ast ((b \ast x)x^{-1})=(b \ast x)x^{-1},
              \end{equation}
              for any \(b\), \(b' \in B\) and \(x \in X\). In this case, the image \(\varphi(B)\) is an abelian subgroup of \(\Aut_c(X)\).
    \end{itemize}
\end{theorem}

\begin{proof}{\ }
    \begin{itemize}
        \item[(i)] We define \(\rho\) as follows: for every \(2\)-nilpotent group \(B\), the component
              \[
                  \rho_B \colon \SplExt(B,X) \to \Hom_\Grp(U(B),\Aut_c(X))
              \]
              is the morphism in \(\Set\) which maps any action
              \[
                  \xi \colon B \times X \to X \colon (b,x) \mapsto b \ast x
              \]
              to the group homomorphism
              \[
                  \varphi \colon B \to \Aut_c(X) \colon b \mapsto b \ast -.
              \]
              The transformation \(\rho\) is natural. Indeed, let
              \[
                  f \colon B' \to B
              \]
              be a morphism in \(\Nil_2(\Grp)\), and let
              \[
                  f^*=\SplExt(f,X) \colon \SplExt(B,X) \to \SplExt(B',X)
              \]
              be the \emph{change of base} function given by pulling back along \(f\). Then, it is easy to verify that the diagram in \(\Set\)
              \[
                  \begin{tikzcd}
                      \SplExt(B,X) \arrow[r, "\rho_B"] \arrow[d, swap, "{f^\ast}"]
                      & \Hom_\Grp(U(B),\Aut_c(X)) \arrow[d,"{- \circ f}"]\\
                      \SplExt(B',X) \arrow[r, swap, "\rho_{B'}"] & \Hom_\Grp(U(B'),\Aut_c(X))
                  \end{tikzcd}
              \]
              is commutative. Moreover, the map \(\rho_B\) is an injection since every element of \(\SplExt(B,X)\) is uniquely determined by the corresponding derived action of \(B\) on \(X\), that is, by the map
              \[
                  \xi \colon B \times X \to X \colon (b,x) \mapsto b \ast x.
              \]
              Therefore, \(\rho\) is a natural monomorphism of functors.
        \item[(ii)] By \Cref{rmk}, a morphism \(\varphi \colon U(B) \to \Aut_c(X)\) defines an action of \(B\) on \(X\) in \(\Nil_2(\Grp)\) if and only if \eqref{act_mor} holds. Moreover, as observed in the proof of \Cref{PropAct}, the condition \eqref{act_mor} implies that
              \[
                  (bb') \ast x = (b'b) \ast x,
              \]
              for any \(b\), \(b' \in B\) and \(x \in X\). This means that the automorphisms \(b \ast -\) and \(b' \ast -\) commute, that is, \(\Imm(\varphi)\) is an abelian subgroup of \(\Aut_c(X)\).\qedhere
    \end{itemize}
\end{proof}

We conclude this section with an example of a canonical action of any \(2\)-nilpotent group \(X\) on itself given by inner automorphisms. Before doing so, we state the following lemma, whose proof is immediate.

\begin{lemma}\label{lemma}

    Let \(X\) be a group, and let \(\Inn(X)\) denote the group of \emph{inner automorphisms} of \(X\). Then
    \[
        \Inn(X) \leq \Aut_c(X)
    \]
    if and only if \(X\) is \(2\)-nilpotent. \noproof
\end{lemma}

\begin{example}
    Let \(X\) be a \(2\)-nilpotent group. Then, by \Cref{lemma}, \(\Inn(X)\) is a subgroup of \(\Aut_c(X)\). %Indeed, let \(g \in X\) and denote by \(\varphi_g\)  the inner automorphism defined by
    %\[
    %\varphi_g(x)=gxg^{-1},
    %\]
    %for every \(x \in X\). Since \([X,X] \leq \bZ(X)\), we have that
    %\[
    %\varphi_g(x)x^{-1}=gxg^{-1}x^{-1}=[g,x]\in \bZ(X)
    %\]
    %and \(\varphi_g\) is a central automorphism of \(X\). Thus
    %\[
    %\Inn(X)\leq \Aut_c(X).
    %\]
    Now consider the canonical group homomorphism
    \[
        \Inn \colon X \to \Inn(X) \colon g \mapsto \varphi_g,
    \]
    where \(\varphi_g\)  is the inner automorphism defined by
    \[
        \varphi_g(x)=gxg^{-1},
    \]
    for any \(x \in X\). We have
    \[
        \varphi_h(\varphi_g(x)x^{-1})=\varphi_h([g,x])=[g,x]=\varphi_g(x)x^{-1},
    \]
    since \([g,x] \in \bZ(X)\). Therefore, by (2) of \Cref{main_thm}, if \(i \colon \Inn(X) \hookrightarrow \Aut_c(X)\) denotes the inclusion, then
    \[
        \overline{\Inn} = i \circ \Inn \in \Imm(\rho_X).
    \]
    Finally, the action of \(X\) on itself induced by \(\overline{\Inn}\) is precisely the conjugation action:
    \[
        g \ast x = \varphi_g(x)=gxg^{-1},
    \]
    for any \(g\), \(x \in X\). It follows that the corresponding semidirect product \(X \ltimes X\) is \(2\)-nilpotent.
\end{example}

\section{Weak action representability of \texorpdfstring{\(2\)}{2}-nilpotent groups}\label{sec_rep}

The result of \Cref{main_thm} may be viewed as a particular case of~\cite[Corollary~4.2]{CigoliManciniMetere}. Indeed, \(\Aut_c(X)\) is a universal strict general actor of \(X\) in the Orzech category of interest \(\Nil_2(\Grp)\). However, since \(\Aut_c(X)\) is not, in general, \(2\)-nilpotent, it does not provide a weak representing object for the actions on \(X\) in the variety \(\Nil_2(\Grp)\). In this section, we show that it is nevertheless possible to construct such a weak representing object starting from the images of the morphisms
\[
    U(B) \to \Aut_c(X)
\]
which correspond to the actions of \(B\) on \(X\) in \(\Nil_2(\Grp)\). Our aim is therefore to prove that \(\Nil_2(\Grp)\) is a weakly action representable category. We start by observing the following.

\begin{proposition}\label{thm_ar}
    The category \(\Nil_2(\Grp)\) is action accessible and not action representable.
\end{proposition}

\begin{proof}
    The category \(\Nil_2(\Grp)\) is action accessible, since it is an Orzech category of interest (see~\cite{Montoli}). On the other hand, by~\cite[Theorem 3.6]{Casas}, it is not action representable, because \(\USGA(X)=\Aut_c(X)\) is not, in general, an object of~\(\Nil_2(\Grp)\).
\end{proof}

We now explain how the problem of finding a weak representation of the actions on a \(2\)-nilpotent group \(X\) is related to the so-called \emph{amalgamation property} (AP)~\cite{kiss}.

We recall that for a diagram in a category \(\C\), an \emph{amalgam} is a monic cocone, i.e., a cocone which is a monomorphic natural transformation. This means that each component of that cocone is a monomorphism, which implies that all the morphisms of the given diagram were monomorphisms to begin with. Note that in a category with colimits, an amalgam for a diagram exists if and only if its colimit cocone is such an amalgam. A category is said to have the amalgamation property~(AP) when every span of monomorphisms admits an amalgam; equivalently, for each pushout square
\[
    \xymatrix{I \ar[r]^t \ar[d]_s & T\ar[d]^-{\iota_T}\\
    S \ar[r]_-{\iota_S} & S+_IT}
\]
if \(s\) and \(t\) are monomorphisms then so are \(\iota_S\) and \(\iota_T\). It is well known that the categories of groups, abelian groups and Lie algebras satisfy the condition (AP), whereas, for instance, the category of associative algebras does not. We refer the reader to~\cite{kiss} for an overview of examples and references.

The relationship between action representability and (AP) was first investigated in~\cite{BJK2}, where it was proved that, for any Orzech category of interest, action representability is equivalent to (AP) for \emph{protosplit} monomorphisms, i.e., monomorphisms arising as kernels of split extensions. Moreover, in~\cite{WRAAlg}, the authors showed that the existence of a weak representation of the actions on a non-associative algebra \(X\) is related to the existence of an amalgam for a certain diagram of subalgebras of the external weak actor \(\E(X)\).

Now, by \Cref{main_thm} each action \(\xi\) of \(B\) on \(X\) in \(\Nil_2(\Grp)\) gives rise to a morphism \(\rho_B(\xi)\colon U(B)\to \Aut_c(X)\) in \(\Grp\). We want to show that \(\xi\) is also determined by a morphism of \(2\)-nilpotent groups \(\tau_B(\xi)\colon B\to T\), where \(T=T(X)\) is a weak representing object of \(X\).

Each morphism \(\rho_B(\xi)\colon U(B)\to \Aut_c(X)\) is the composite of a surjective morphism of \(2\)-nilpotent groups \(\rho'_B(\xi)\colon U(B)\twoheadrightarrow M_\xi\), where \(M_\xi=\Imm(\rho_B(\xi))\), and the canonical inclusion \(M_\xi \hookrightarrow \Aut_c(X)\). We thus find a diagram of \(2\)-nilpotent subgroups of \(\Aut_c(X)\) indexed over the \(2\)-nilpotent group actions on \(X\). We may re-index and view \(\{M_\xi\}_\xi\) as a diagram in \(\Grp\) over a subcategory of the thin category of \(2\)-nilpotent subgroups of \(\Aut_c(X)\). We further observe that, by \Cref{PropAct}, each \(M_\xi\) is an abelian group. Thus, \(\{M_\xi\}_\xi\) is a family of abelian subgroups of the group of central automorphisms of~\(X\).

Now, the category \(\Ab\) of abelian groups has the amalgamation property~\cite{kiss}, and we can therefore construct an amalgam\footnote{We recall that, given a span of monomorphisms \(B \mathrel{\stackrel{f}{\leftarrowtail}} S \mathrel{\stackrel{f'}{\rightarrowtail}} B'\) in \(\Ab\), its amalgam is given by the pushout
    \(B +_S B'=(B \oplus B') / H\), where \(H=\langle (f(s),f'(s)^{-1}) \mid s \in S \rangle\).} \(T\) of \(\{M_\xi\}_\xi\) in \(\Ab\). Thus, we can associate with any action~\(\xi\) of \(B\) on \(X\) in the variety \(\Nil_2(\Grp)\), the homomorphism
\begin{equation*}
    \begin{tikzcd}
        B \arrow[r, two heads, "\rho'_B(\xi)"]
        & M_\xi \arrow[r, tail, "j_{M_\xi}"]
        & T
        \arrow[from=1-1, to=1-3, bend right=25, "\tau_B(\xi)"']
    \end{tikzcd}
\end{equation*}
where \(j_{M_\xi}\) denotes the inclusion of \(M_\xi\) inside the amalgam \(T\). We further observe that \(j_{M_\xi}\) depends only on the abelian group \(M_\xi\) and not on the action \(\xi\). We now prove that the assignment \(\xi \mapsto \tau_B(\xi)\) defines a weak representation of the actions on \(X\), as stated in the following theorem.

\begin{theorem}\label{thm_wra}
    The category \(\Nil_2(\Grp)\) is weakly action representable. Moreover, for any \(2\)-nilpotent group \(X\), a weak representing object of \(X\) is the abelian group~\(T\) constructed above.
\end{theorem}

\begin{proof}
    Let \(X\) be a \(2\)-nilpotent group. By the discussion above, for any other \(2\)-nilpotent group \(B\), we define a map
    \[
        \tau_B \colon \SplExt(B,X) \to \Hom_{\Nil_2(\Grp)}(B,T).
    \]
    Such a map is injective. Indeed, suppose there exist two actions \(\xi\) and \(\psi\) of \(B\) on~\(X\) in the variety \(\Nil_2(\Grp)\) such that \(\tau_B(\xi)=\tau_B(\psi)\). Then, by uniqueness of image factorisations, the images of \(j_{M_\xi}\colon M_\xi\to T\) and \(j_{M_\psi}\colon M_\psi\to T\) are isomorphic subgroups of \(T\). Now, the image in \(\Grp\) of the diagram \(\{M_\xi\}_\xi\) is a thin category, so that \(M_\xi=M_\psi\). Hence
    \[
        j_{M_\xi}\circ \rho'_B(\xi)=\tau_B(\xi)=\tau_B(\psi)=j_{M_\psi}\circ\rho'_B(\psi)=j_{M_\xi}\circ\rho'_B(\psi),
    \]
    which implies \(\rho'_B(\xi)=\rho'_B(\psi)\). Since \(\rho\) is a natural monomorphism, we get \(\xi=\psi\).

    It remains to prove that the collection \(\{\tau_B\}_B\) gives rise to a natural transformation
    \[
        \tau \colon \SplExt(-,X) \to \Hom_{\Nil_2(\Grp)}(-,T).
    \]
    This follows directly by the naturality of both \(\rho'\) and the cocone components in the amalgam. In fact, if two maps \(\rho_B(\xi)\) and \(\rho_C(\psi)\) defined from \(B\) and \(C\), respectively, to \(\Aut_c(X)\) have the same image subgroup \(M_\xi=M_\psi\) of \(\Aut_c(X)\), then by naturality of \(\rho\), and the fact that the inclusion of \(M_\xi\) into \(\Aut_c(X)\) is a monomorphism, for any equivariant morphism \(f \colon B\to C\) we have that the square on the left in
    \[
        \xymatrix{B \ar[d]_-{f} \ar[r]^-{\rho'_B(\xi)} & M_\xi \ar@{=}[d] \ar[r]^-{j_{M_\xi}} & T \ar@{=}[d]\\
        C \ar[r]_-{\rho'_C(\psi)}  & M_\psi \ar[r]_-{j_{M_\psi}} & T}
    \]
    commutes. Since the entire diagram is commutative, we conclude that \(\tau\) is a natural monomorphism of functors. Thus, the category \(\Nil_2(\Grp)\) is weakly action representable, and a weak representing object of \(X\) is the abelian group \(T\).
\end{proof}

The result of \Cref{thm_wra} settles an open question left in~\cite{XabiMancini}, where it is shown that, for \(k \geq 3\), the category \(\Nil_k(\Grp)\) is not weakly action representable, and at the same time establishes a connection with the case of \(2\)-nilpotent Lie algebras (see~\cite[Theorem 2.21]{WRAAlg}, where it is proved that the variety of \(2\)-nilpotent anti-commutative algebras is weakly action representable).

As explained in~\cite[Corollary~5.3 and Corollary 5.4]{WRA}, the existence of the initial weak representation follows from the existence of a weak representation and the fact that \(\Nil_2(\Grp)\) is a total category~\cite{Street}.


\section{A note on local algebraic cartesian closedness}\label{sec_LACC}

We conclude the article with a proof that the variety \(\Nil_2(\Grp)\) is not \emph{locally algebraically cartesian closed} \LACC{}. As explained in the introduction, this condition was introduced by J.~R.~A.~Gray in~\cite{GrayThesis} in connection with the concept of \emph{algebraic exponentiation}, and it is closely related to the study of internal actions in semi-abelian categories.

Recall that a semi-abelian category \(\C\) is said to be \emph{locally algebraically cartesian closed} when the change-of-base functors in the fibration of points are left adjoints. This is equivalent to the condition that for every object \(B\) of \(\C\), the functor
\[
    B\flat(-)\colon \C \to \C,
\]
where for any object \(X\) the object \(B\flat X\) is defined as the kernel of the canonical morphism
\[
    [\id_B,0]\colon B+X\to B,
\]
is a left adjoint. As explained in~\cite{GrayACS}, when \(\C\) is a variety of algebras, this is further equivalent to the conditions
\begin{enumerate}
    \item \(B\flat(-)\) preserves all colimits;
    \item \(B\flat(-)\) preserves binary sums.
\end{enumerate}
Hence, \(\C\) is \LACC{} if and only if, for every pair of objects \(X\) and \(Y\) of \(\C\), the canonical comparison morphism
\begin{equation}\label{comparison}
    [B \flat \id_X, B \flat \id_Y] \colon (B\flat X)+(B\flat Y)\to B\flat(X+Y)
\end{equation}
is an isomorphism.

Examples of semi-abelian categories satisfying the condition \LACC{} include the category of groups, the variety of Lie algebras over a unital commutative ring~\cite{GrayLie}, and the category of cocommutative Hopf algebras over a field of characteristic zero, see~\cite{acc, hopf}. It was then proved in~\cite{GM-VdL3} that, among varieties of non-associative algebras over an infinite field of characteristic different from \(2\), the only non-abelian \LACC{} variety is the category of Lie algebras. Thus, for varieties of non-associative algebras, action representability is equivalent to the condition \LACC{}. On the other hand, for arbitrary semi-abelian categories this is known to fail in general. Indeed, the category of Boolean rings is action representable~\cite{BJK2} but does not satisfy \LACC{}, as explained in~\cite[Proposition 6.4]{GrayACS}. It remains an open problem whether local algebraic cartesian closedness implies action representability, or whether a counterexample exists. This makes us wonder what happens in the case of \(2\)-nilpotent groups. We prove the following result:

\begin{theorem}
    The variety \(\Nil_2(\Grp)\) is not locally algebraically cartesian closed.
\end{theorem}

\begin{proof}
    It is sufficient to find \(2\)-nilpotent groups \(B\), \(X\), and \(Y\) such that \eqref{comparison} is not a monomorphism. Here, when we write \(G=\langle S \mid R \rangle_{\Nil_2}\), we mean that \(G\) is the group generated by \(S\)
    subject to the relations \(R\) in the variety \(\Nil_2(\Grp)\). Equivalently, \(G\) is the quotient of the free \(2\)-nilpotent group on \(S\) by the normal subgroup generated by \(R\). In
    particular, all commutators are tacitly central.

    Let \(B=\langle b\rangle\), \(X=\langle x\rangle\), and \(Y=\langle y\rangle\) be three copies of the infinite cyclic group~\(\Z\). One may check that
    \begin{align*}
        B \flat X     & = \langle x,u \mid [x,u]=1 \rangle_{\Nil_2} \cong \Z \times \Z,                           \\
        B \flat Y     & = \langle y,v \mid [y,v]=1 \rangle_{\Nil_2} \cong \Z \times \Z,                           \\
        B \flat (X+Y) & =\langle x,y,u,v \mid [x,u]=[y,v]=[u,v]=1 \rangle_{\Nil_2} \cong H_3 \times \Z \times \Z,
    \end{align*}
    where \(u=[b,x]\), \(v=[b,y]\), and \(H_3\) denotes the discrete Heisenberg group.

    Furthermore, the coproduct \((B \flat X) + (B \flat Y)\) is the group
    \[
        \langle x,y,u,v \mid [x,u]=[y,v]=1 \rangle_{\Nil_2},
    \]
    whose commutator subgroup is the abelian group generated by
    \[
        [x,y], \quad [x,v], \quad [y,u], \quad [u,v].
    \]
    We remark that none of these four commutators is trivial.

    Now consider the comparison map
    \[
        \varphi=[B \flat \id_X,\, B \flat \id_Y]\colon (B\flat X)+(B\flat Y)\to B\flat(X+Y).
    \]
    It is defined on generators by
    \[
        x \mapsto x,\quad y \mapsto y,\quad u \mapsto u,\quad v \mapsto v,
    \]
    where we use the same symbols for generators in the domain and codomain.

    Since \(u\) and \(v\) are central in \(B\flat (X+Y)\), it follows that
    \[
        \varphi([x,v])=\varphi([y,u])=\varphi([u,v])=1.
    \]
    Hence \(\Ker(\varphi)\) is non-trivial, so \(\varphi\) is not an isomorphism. Therefore, the category \(\Nil_2(\Grp)\) is not \LACC{}.
\end{proof}

This result suggests a possible direction for future research, namely to investigate whether the failure of the condition \LACC{} extends more generally to arbitrary non-abelian subvarieties of the category \(\Grp\).

\section*{Acknowledgements}
The first and second authors are supported by the University of Palermo and by the ``National Group for Algebraic and Geometric Structures and their Applications'' (GNSAGA -- INdAM).

The second author is also supported by the SDF Sustainability Decision Framework Research Project -- MISE decree of 31/12/2021 (MIMIT Dipartimento per le politiche per le imprese -- Direzione generale per gli incentivi alle imprese) -- CUP:~B79J23000530005, COR:~14019279, Lead Partner:~TD Group Italia Srl, Partner:~University of Palermo, and he is a Postdoctoral Researcher of the Fonds de la Recherche Scientifique--FNRS.

The third author is a Senior Research Associate of the Fonds de la Recherche Scientifique--FNRS.

%\section*{Declarations}
%Not applicable. There is no Competing Interest.

\begin{thebibliography}{99}

    \bibitem{Barr}
    M.~Barr, \emph{Exact Categories}, Lecture Notes in Mathematics \textbf{236} (1971), 1--120.

    \bibitem{BJK}
    F.~Borceux, G.~Janelidze and G.~M. Kelly, \emph{Internal object actions}, Commentationes Mathematicae Universitatis Carolinae \textbf{46} (2005), no.~2, 235--255.

    \bibitem{BJK2}
    F.~Borceux, G.~Janelidze and G.~M. Kelly, \emph{On the representability of actions in a semi-abelian category}, Theory and Applications of Categories \textbf{14} (2005), no.~11, 244--286.

    \bibitem{Bourn}
    D.~Bourn, \emph{Normalization equivalence, kernel equivalence and affine categories}, Lecture Notes in Mathematics \textbf{1488} (1991), 43--62.

    \bibitem{BournGray}
    D.~Bourn and J.~R.~A.~Gray, \emph{Aspects of algebraic exponentiation}, Bulletin of the Belgian Mathematical Society -- Simon Stevin \textbf{19} (2012), no.~5, 823--846. \doi{10.36045/bbms/1354031552}

    \bibitem{act_accessible}
    D.~Bourn and G.~Janelidze, \emph{Centralizers in action accessible categories}, Cahiers de Topologie et G\'eom\'etrie Diff\'erentielle Cat\'egoriques \textbf{50} (2009), no.~3, 211--232.

    \bibitem{WRAAlg}
    J.~Brox, X.~García-Martínez, M.~Mancini, T.~Van~der Linden and C.~Vienne, \emph{Weak representability of actions of non-associative algebras}, Journal of Algebra \textbf{669} (2025), no.~18, 401--444. \doi{10.1016/j.jalgebra.2025.02.007}

    \bibitem{Casas}
    J.~M.~Casas, T.~Datuashvili and M.~Ladra, \emph{Universal strict general actors and actors in categories of interest}, Applied Categorical Structures \textbf{18} (2010), 85--114. \doi{10.1007/s10485-008-9166-z}

    \bibitem{acc}
    A.~S. Cigoli, J.~R.~A. Gray and T.~Van~der Linden, \emph{Algebraically coherent categories}, Theory and Applications of Categories \textbf{30} (2015), no.~54, 1864--1905.

    \bibitem{CigoliManciniMetere}
    A.~S.~Cigoli, M.~Mancini and G.~Metere, \emph{On the representability of actions of Leibniz algebras and Poisson algebras}, Proceedings of the Edinburgh Mathematical Society \textbf{66} (2023), no.~4, 998--1021. \doi{10.1017/S0013091523000548}

    \bibitem{Datuashvili}
    T. Datuashvili, \emph{Cohomologically trivial internal categories in categories of groups with operations}, Applied Categorical Structures \textbf{3} (1995), 221--237. \doi{10.1007/BF00878442}

    \bibitem{XabiMancini}
    X.~García-Martínez and M.~Mancini, \emph{Action accessible and weakly action representable varieties of algebras}, Bollettino dell'Unione Matematica Italiana (2025). \doi{10.1007/s40574-025-00496-1}

    \bibitem{GM-VdL2}
    X.~García-Martínez and T.~Van der Linden, \emph{A characterisation of Lie algebras amongst anti-commutative algebras}, Journal of Pure and Applied Algebra \textbf{223} (2019), no.~11, 4857--4870.  \doi{10.1016/j.jpaa.2019.02.018}

    \bibitem{GM-VdL3}
    X.~García-Martínez and T.~Van der Linden, \emph{A characterisation of Lie algebras via algebraic exponentiation}, Advances in Mathematics \textbf{341} (2019), 92--117. \doi{10.1016/j.aim.2018.10.034}

    \bibitem{Tim}
    X.~García-Martínez, M.~Tsishyn, T.~Van der Linden and C.~Vienne, \emph{Algebras with representable representations}, Proceedings of the Edinburgh Mathematical Society \textbf{64} (2021), no.~2, 555--573. \doi{10.1017/S0013091521000304}

    \bibitem{hopf}
    M.~Gran, G.~Kadjo, and J.~Vercruysse, \emph{A torsion theory in the category of cocommutative Hopf algebras}, Applied Categorical Structures \textbf{24} (2016), 269--282. \doi{10.1007/s10485-015-9396-9}

    \bibitem{GrayLie}
    J.R.A.~Gray, \emph{Algebraic exponentiation for categories of Lie algebras}, Journal of Pure and Applied Algebra \textbf{216} (2012), no.~8--9, 1964--1967. \doi{10.1016/j.jpaa.2012.02.034}

    \bibitem{GrayThesis}
    J.~R.~A. Gray, \emph{Algebraic Exponentiation in General Categories}, Ph.D. thesis, University of Cape Town, 2010.

    \bibitem{GrayACS}
    J.~R.~A.~Gray, \emph{Algebraic exponentiation in general categories}, Applied Categorical Structures \textbf{20} (2012), no.~6, 543--567. \doi{10.1007/s10485-011-9251-6}

    \bibitem{Gray}
    J.~R.~A. Gray, \emph{A note on the relationship between action accessible and weakly action representable categories}, Theory and Applications of Categories \textbf{44} (2025), no.~8, 272--276.

    \bibitem{WRA}
    G.~Janelidze, \emph{Central extensions of associative algebras and weakly action representable categories}, Theory and Applications of Categories \textbf{38} (2022), no.~36, 1395--1408.

    \bibitem{kiss}
    E.~W.~Kiss, L.~Márki, P.~Pröhle and W.~Tholen, \emph{Categorical algebraic properties. A compendium on amalgamation, congruence extension, epimorphisms, residual smallness, and injectivity}, Studia Scientiarum Mathematicarum Hungarica \textbf{18} (1983), 79--141.

    \bibitem{Semi-Ab}
    G.~Janelidze, L.~Márki and W.~Tholen, \emph{Semi-abelian categories}, Journal of Pure and Applied Algebra \textbf{168} (2002), no.~2, 367--386. \doi{10.1016/S0022-4049(01)00103-7}

    \bibitem{loday}
    J.-L.~Loday, \emph{Une version non commutative des algèbres de Lie: les algèbres de Leibniz}, L'Enseignement Math\'ematique \textbf{39} (1993), no.~3-4, 269--293.

    \bibitem{Mancini_Thesis}
    M.~Mancini, \emph{On the representability of actions of non-associative algebras}, Ph.D. Thesis (2024), University of Palermo, Italy. \url{https://hdl.handle.net/10447/623093}

    \bibitem{Montoli}
    A.~Montoli, \emph{Action accessibility for categories of interest}, Theory and Applications of Categories \textbf{23} (2010), no.~1, 7--21.

    \bibitem{Neumann}
    B.~H.~Neumann, \emph{Permutational products of groups}, Journal of the Australian Mathematical Society \textbf{1} (1960), no.~3, 299--310.

    \bibitem{Orzech}
    G.~Orzech, \emph{Obstruction theory in algebraic categories {I} and {II}}, Journal of Pure and Applied Algebra~\textbf{2} (1972), no.~4, 287--314 and 315--340. \doi{10.1016/0022-4049(72)90008-4}

    \bibitem{Street}
    R.~Street and R.~Walters, \emph{Yoneda structures on 2-categories}, Journal of Algebra \textbf{50} (1978), no.~2, 350--379. \doi{10.1016/0021-8693(78)90160-6}

    \bibitem{VdL-NAA}
    T.~Van~der Linden, \emph{Non-associative algebras}, New Perspectives in Algebra, Topology and Categories (M.~M. Clementino, A.~Facchini, M.~Gran, eds.), Coimbra Mathematical Texts, vol.~1, Springer, Cham, 2021, 225--258. \doi{\detokenize{10.1007/978-3-030-84319-9_7}}
\end{thebibliography}
\end{document}
